% Generated from 02-具身智能的问题定义、系统边界与核心矛盾.md; edit the LaTeX source after migration.

\part*{第二部分 具身智能的问题定义、系统边界与核心矛盾}
\addcontentsline{toc}{part}{第二部分 具身智能的问题定义、系统边界与核心矛盾}
\markboth{第二部分 具身智能的问题定义、系统边界与核心矛盾}{第二部分 具身智能的问题定义、系统边界与核心矛盾}

\chapter*{6. 什么是具身智能}
\addcontentsline{toc}{chapter}{6. 什么是具身智能}
\markboth{6. 什么是具身智能}{6. 什么是具身智能}

为了让这一定义在工程上真正可用，还需要加上一条常被忽略的条件：具身智能讨论的不是“是否具有身体”这一静态事实，而是“身体是否构成问题结构的一部分”。一个相机支架也有物理形体，但如果它不承担自主感知、任务选择和动作后果，那么它并不进入本报告意义上的具身系统。相反，一个强依赖远程操作者、但仍须在局部时间尺度上独立完成接触控制和状态维持的遥操作平台，则已经部分落入具身系统的讨论范围。也就是说，“具身”不是二元标签，而更像一个系统责任划分问题。

“具身智能”最直观的含义，是智能不再仅仅表现为符号处理、文本生成或抽象决策，而是必须通过一个具有物理身体、传感通道和执行能力的系统，在真实或近真实环境中完成感知、判断、行动和反馈闭环。本部分的职责，不只是给出一个定义，而是把“具身智能”重新压回一个可操作的问题定义之中：哪些系统应被纳入讨论，哪些相关概念必须区分，后文将用什么分层框架分析系统，以及为什么当前这个领域会被几组核心矛盾持续牵引。换句话说，只有当一个系统在与环境持续耦合的过程中展现出稳定、可迁移、可纠偏的行为能力时，我们才有充分理由把它视为具身的智能系统，而不只是一个外挂到机器人上的计算模块。

这一概念之所以重要，是因为它把“智能”的评价标准从离线任务成绩，转向了动态环境中的行为表现。在文本任务里，一个模型回答错误，通常只是输出一段不准确的文字；在具身系统里，一个错误的感知、规划或动作，可能导致抓取失败、跌倒、碰撞、误伤、任务中断，甚至造成设备损坏。因此，具身智能不是简单把语言模型装进机器人，而是要求整个系统在物理世界约束下建立持续有效的闭环能力。

如果进一步拆解，“具身”至少包含三层含义。第一，系统具有载体，不论该载体是固定机械臂、移动平台、四足、人形还是无人机，它都必须承担真实物理执行。第二，系统与环境之间存在持续的信息流和作用流，也就是传感与控制的双向闭环。第三，系统行为必须受到时间、空间、接触、能耗、安全和任务目标等约束。这三层含义共同决定，具身智能讨论的并不是纯认知问题，而是“认知如何嵌入物理行动”这一系统问题。

因此，本报告中的“具身智能”不是一个营销口号，而是一种对问题结构的界定：一个系统是否具身，不取决于它是否使用了大模型，也不取决于它是否长得像人，而取决于它是否必须在物理世界中以闭环方式完成任务，并为自己的行动后果承担实时约束。

如果用论证结构来表达，本章在这里提出的核心 claim 是：具身智能首先是一个“系统问题定义”，而不是“模型类别名称”。支持这一 claim 的 data 来自三个方面。其一，物理载体与环境耦合决定了系统必须面对接触、延迟、噪声、失稳、失效恢复和安全等问题，这些问题在纯软件系统中要么不存在，要么代价完全不同。其二，真实任务的完成依赖感知、表征、规划、控制与执行之间的闭环协同，而不是单点推理能力。其三，当前行业中真正困难和昂贵的部分，往往并不出现在“理解一句指令”上，而出现在“把理解稳定转成可持续物理行为”上。

这一工作定义还有一个重要作用，是压低概念膨胀的风险。因为“具身智能”在公共叙事中很容易被同时拿来指代机器人学的新名字、类人机器人的愿景、大模型的物理延伸，甚至更宏大的通用智能未来。如果报告不先把这个概念重新压回“物理世界中的闭环行为系统”这一最小可操作定义，那么后文的比较对象就会不断变化，最后难以形成稳定判断。

若把“具身系统”进一步写成一个最小状态闭环，它至少可以被抽象为：

\begingroup
\footnotesize
\begin{equation*}
\mathcal{E} = (\mathcal{B},\ \mathcal{O},\ \mathcal{A},\ \mathcal{M},\ \Pi,\ \mathcal{C})
\end{equation*}
\endgroup

其中 \(\mathcal{B}\) 是本体与执行器约束，\(\mathcal{O}\) 是观测通道，\(\mathcal{A}\) 是动作空间，\(\mathcal{M}\) 是内部状态或表征，\(\Pi\) 是策略或控制接口，\(\mathcal{C}\) 则是时间、能耗、安全、接触与任务规则等约束集合。这个写法的重要性在于，它迫使我们把“身体”“感知”“动作”“内部模型”和“物理约束”同时作为问题结构的一部分，而不是只盯住其中最显眼的模型层。

\chapter*{7. 具身智能与相关概念的关系}
\addcontentsline{toc}{chapter}{7. 具身智能与相关概念的关系}
\markboth{7. 具身智能与相关概念的关系}{7. 具身智能与相关概念的关系}

这里还应补充一个阅读上的提醒：这些概念之间不是互斥集合，而更像若干相互穿插的投影面。自动驾驶、移动操作、家庭机器人、工业协作臂、人机协作系统都可能同时共享部分具身能力结构；区别不在于是否“属于具身智能阵营”，而在于它们分别把难度压在哪些层上。把这些对象强行放进单一概念桶里，往往会遮蔽真正重要的比较维度。

具身智能并不是凭空出现的新学科，而是建立在机器人学长期积累之上的一次重新组织。传统机器人学关心的运动学、动力学、规划、控制、定位、建图、抓取与系统集成，本质上都在回答“如何让机器在物理世界中行动”这一问题。具身智能之所以看起来像一个新范式，不是因为这些旧问题失效了，而是因为大模型、多模态学习和大规模数据驱动方法，把原本相对分散的能力模块重新拉回到了一个更统一的表征与决策框架之下。

它与传统自动化/控制之间既连续又不同。连续性在于，两者都处理真实物理系统，都要面对动态约束、误差累积与安全问题；差异在于，传统自动化通常工作在更强结构化、更弱语义需求的环境中，而具身智能更关心开放场景中的任务理解、环境泛化和多模态决策。换言之，具身智能不是“不要控制了”，而是在控制之上叠加了更复杂的感知与认知要求。

它与移动机器人、自动驾驶也存在重叠。移动机器人和自动驾驶本来就是具身系统的重要组成部分，因为它们同样具有身体、感知、规划和执行闭环。但自动驾驶的任务边界、监管约束、地图依赖和安全要求具有显著特化，因而不能简单代表“通用具身智能”的全部问题。自动驾驶更像是一个规则密集、场景复杂但任务定义相对集中化的具身子领域，而通用机器人则常常面临更加多样化的对象、任务和交互形式。

它与大模型智能体的关系同样需要厘清。纯软件智能体可以具备任务分解、工具调用、外部记忆和多步推理能力，但如果它不直接承担物理执行约束，那么它更适合作为具身系统的高层参考模块，而不能直接等同于具身系统本身。具身智能可以借用大模型智能体的高层能力，但必须把这些能力重新映射到实时控制、接触行为和环境风险可控的执行链路中。

最后，它与 AGI 叙事也有交集但不应混同。具身智能经常被纳入“通用智能走向真实世界”的宏大叙事中，但本报告会尽量避免把“具身智能”直接写成 AGI 的线性前置阶段。原因在于，具身智能首先是一个工程与系统问题，它的难点很大程度上来自物理约束、数据稀缺、部署复杂性和安全问题，而不是抽象推理能力本身。

本节最重要的意义，不是给出一张看似全面的概念关系图，而是建立一套足以支持后续比较的边界纪律。若不先厘清具身智能与传统自动化、移动机器人、自动驾驶、软件智能体和 AGI 叙事之间的关系，后文所有关于“哪些路线更重要”“哪些企业更值得跟踪”“哪些论文在解决核心问题”的判断都会失去统一标准。边界工作看似是导论工作，实则决定了后文比较能否成立。

更严格地说，本节真正想建立的是一套“比较坐标”而不是“概念名册”。对后文更有用的问题并不是“某对象到底算不算具身智能”，而是“它主要在哪些层上承担责任”。例如，自动驾驶在环境感知、状态估计和规则约束上高度成熟，但在开放对象操作和精细接触上并不能直接代表通用机器人；软件智能体在高层规划、工具调用与外部记忆上可作为参考，但若不承担物理闭环与动作后果，就不应与具身系统等同比较。把对象按责任层而不是按营销标签归类，后文才不容易发生口径漂移。

\chapter*{8. 具身智能系统的分析框架}
\addcontentsline{toc}{chapter}{8. 具身智能系统的分析框架}
\markboth{8. 具身智能系统的分析框架}{8. 具身智能系统的分析框架}

这套分层还隐含着一个很重要的写作原则：后文讨论任何“突破”时，都应尽量回答它究竟改善了哪一层、通过什么接口影响相邻层、又在哪些层面仍保留旧瓶颈。只有这样，报告才不会把局部提升误写成全栈突破。例如，一个新 VLA 可能显著提升任务描述理解与动作先验，却未必改写接触建模和低层稳定控制；一个强世界模型可能改善表征与预期评估，却未必直接解决真实部署中的计算预算和安全验证问题。

为了避免把具身智能误解为某个单点模型，本报告采用一个分层分析框架来组织后续内容。第一层是本体层，关心机器人的机械结构、自由度、执行器、末端执行器、传感器布局、算力和供能约束。第二层是感知层，关心系统如何从视觉、触觉、力觉、本体感觉与状态信号中构造对环境和自身状态的估计。第三层是表征与世界建模层，关心如何把离散观测组织成可用于预测、推理和规划的内部表示。第四层是决策与规划层，关心如何把任务目标转化为子任务、策略和行动计划。第五层是控制与执行层，关心动作如何在实时性、稳定性和安全约束下被真正落实到物理系统上。

这一框架的作用不只是“把问题列成五层”，更重要的是帮助区分不同论文、系统或企业究竟在解决哪一层的问题。许多讨论之所以混乱，原因之一就是不同人谈论的其实不是同一个层级：有人说的是高层语义任务理解，有人说的是动作表征和策略学习，有人说的是低层控制和硬件可靠性，却都使用“具身智能能力强弱”这一单一表述。采用分层框架，可以减少这种概念混淆。

本体层决定了系统的先天能力边界。不同形态、自由度、执行器路线和末端执行器设计，决定了系统能够接触什么对象、进入什么环境、承受多大载荷、实现多高速度以及以何种成本运行。感知层则负责把原始观测转化为对环境和自身状态的稳定估计，这一层直接受传感器类型、布置方式、时间同步和噪声水平影响。表征与世界建模层关心的是“系统在内部如何理解环境”，它决定了系统是否能够从表面观测中形成更稳健、更可迁移的任务相关状态表示。决策与规划层则决定系统如何把任务目标转化为行动结构。控制与执行层最终把这些上层意图落实为安全、稳定、实时的行为。

在这五层之外，还必须有一条横贯所有层级的安全与治理主线。因为具身系统的任何单层失效，都可能累积成真实物理风险。感知错判可能引发错误抓取，世界模型偏差可能导致不合理预期，规划错误可能让系统进入危险状态，控制失稳则可能直接造成物理损伤。因此，本报告后续关于 VLA、世界模型、大小脑分层和端到端路线的讨论，都会建立在这套分层框架上。

采用分层框架还有一个更现实的好处：它能帮助识别许多所谓“技术路线之争”其实是层级错位造成的伪争论。例如，当一条路线声称自己在做“具身基础模型”时，需要进一步追问：它究竟统一的是感知表示、任务理解，还是直接统一到动作生成？当另一条路线强调“大小脑解耦”时，也需要追问：它解决的是高层语义决策与低层实时控制的接口问题，还是仅仅出于工程调度便利？如果没有分层框架，这些路线很容易被放在同一个平面上争论，实际却并不在解决同一层问题。

从系统工程角度看，本章的五层划分并不是唯一正确的理论分法，但它有一个关键优点：它既足够贴近真实机器人系统的工程组织方式，又能容纳大模型时代新增的表征与高层推理能力。因此，这一框架不仅用于解释当前已有系统，也用于作为后文整理新模型、新企业和新架构时的公共挂载点。

若把五层系统接口写成最小递推关系，则可近似表达为：

\begingroup
\footnotesize
\begin{equation*}
o_t \rightarrow z_t \rightarrow g_t \rightarrow u_t \rightarrow x_{t+1}
\end{equation*}
\endgroup

这里 \(o_t\) 是原始观测，\(z_t\) 是任务相关表征，\(g_t\) 是高层目标或阶段性决策，\(u_t\) 是控制量或技能调用，\(x_{t+1}\) 则是下一时刻系统状态。这个链条看似简单，却能帮助读者在后文判断某一论文究竟是在改进哪一段接口，以及其改进是否真的穿透到了后续层，而不是只在局部指标上更好。

\chapter*{9. 具身智能的核心难题}
\addcontentsline{toc}{chapter}{9. 具身智能的核心难题}
\markboth{9. 具身智能的核心难题}{9. 具身智能的核心难题}

若把这五组矛盾视作一种“章节间公共问题清单”，那么它们的价值不在于一句话概括难点，而在于为后文提供复审标准。也就是说，无论后面讨论的是世界模型、扩散策略、OpenVLA、GR00T、灵巧手还是具身公司路线，最终都应能回到这些矛盾上回答：它主要缓解了什么，又把什么问题保留下来甚至放大了。这样全书才不至于在后半段变成平铺直叙的路线罗列。

当前具身智能之所以既令人兴奋，又明显处在早期阶段，是因为它同时受制于几组难以回避的矛盾。第一组矛盾，是高维感知与低时延控制之间的冲突。大模型、多模态模型擅长处理长上下文和复杂语义，但机器人的许多底层控制必须在严格时间预算下运行，这使得统一大模型很难无缝覆盖所有时间尺度。现实系统因此常常不可避免地采用分层设计：高层负责慢速语义决策，低层负责快速稳定控制。

第二组矛盾，是开放世界泛化与安全约束之间的冲突。系统越希望在陌生物体、陌生指令和陌生环境中保持有效行为，就越容易引入不可控失效模式。纯软件系统中的幻觉和误判已经是问题，但在具身系统中，这类错误会被放大为真实物理后果。因此，一个系统是否“更通用”，必须和它是否仍然“可约束、可恢复、可验证”一起被讨论。

第三组矛盾，是数据驱动需求与真实世界数据稀缺之间的冲突。早期视觉直接控制工作 \href{https://arxiv.org/abs/1504.00702}{《End-to-End Training of Deep Visuomotor Policies》} 展示了从视觉直接学习控制策略的可能性，而 \href{https://arxiv.org/abs/1806.10293}{QT-Opt} 则进一步说明，大规模真实机器人交互数据能够显著提升闭环抓取能力；但这类数据的采集成本、设备成本和失败代价都远高于互联网文本或图像数据。具身系统不像语言模型那样能直接从海量廉价公开语料中获益，它必须面对“数据小、场景碎、标注难、失败贵”的现实。

第四组矛盾，是统一端到端与可解释模块化之间的权衡。端到端方法在统一表征和任务泛化上很有吸引力，但调试、验证和安全控制往往更难；模块化系统更可控，但容易牺牲部分通用性与表示共享。今天许多争论，如 VLA 是否应统一到一个超大模型、大小脑应否解耦、动作是否应该离散 token 化，本质上都围绕这一矛盾展开。

第五组矛盾，是研究演示与工程可交付之间的冲突。一个系统可以在论文 benchmark 或实验室演示里取得亮眼结果，但这不自动意味着它能在真实环境中长期、稳定、低维护成本地工作。很多工业化失败并不是败在“没有更强的模型”，而是败在数据采集闭环、硬件寿命、现场运维、安全验证、异常恢复和总拥有成本上。

如果用 Toulmin 式的论证分析来看，本节的每一组“矛盾”都不应被写成绝对断言，而应带有限定词与反驳空间。比如，“统一大模型难以覆盖所有时间尺度”并不是说统一路线一定失败，而是说在当前技术与部署约束下，它通常需要通过分层、技能调用或外部控制回路来弥补时延与稳定性问题。同样，“真实世界数据稀缺”也不是说数据驱动路线必然受阻，而是说其扩展成本显著高于互联网语料型系统。这种带 qualifier 的写法，对于研究型报告很重要，因为它能避免把阶段性工程判断误写成理论定理。

更进一步地说，本节的五组矛盾实际上预先定义了后续章节的主轴。高维感知与低时延控制的矛盾，会在系统分层、动作表征、控制接口和大小脑分工章节中反复出现；开放世界泛化与安全可控的矛盾，会延伸到评测、安全治理与部署约束章节；数据规模与真实世界稀缺性的矛盾，会贯穿数据工程、仿真和训练闭环；统一端到端与模块化可控的矛盾，则几乎是理解当前所有路线分歧的核心入口。

若把这些矛盾压成一个多目标优化问题，则当前具身系统几乎都在求解：

\begingroup
\scriptsize
\begin{equation*}
\min_{\theta}\;
\lambda_1 \mathcal{L}_{\text{semantic}}
+ \lambda_2 \mathcal{L}_{\text{control}}
+ \lambda_3 \mathcal{L}_{\text{safety}}
+ \lambda_4 \mathcal{L}_{\text{data}}
+ \lambda_5 \mathcal{L}_{\text{deployment}}
\end{equation*}
\endgroup

不同路线真正的差别，往往不在于有没有这五项，而在于隐含地把哪一项权重放得更高。VLA 往往更强调语义接口与跨任务共享，经典控制路线更强调低层稳定与安全，数据闭环路线更强调样本覆盖与恢复能力，企业系统则常把部署与运维代价权重拉得更高。把路线差异理解成权重分配差异，比把它们理解成互斥阵营，更接近真实研究与工程世界。

\chapter*{10. 为什么大模型兴起后具身智能再次升温}
\addcontentsline{toc}{chapter}{10. 为什么大模型兴起后具身智能再次升温}
\markboth{10. 为什么大模型兴起后具身智能再次升温}{10. 为什么大模型兴起后具身智能再次升温}

这一判断对后续全书还有一个方法论后果：评价“大模型是否真正改变了机器人”，不能只看是否出现了更会说话、会解释的机器人演示，而应看它是否在任务定义、数据复用、跨场景迁移、异常恢复和部署组织方式上改变了系统结构。如果没有这一层结构变化，那么很多“具身大模型进展”就仍然更接近局部能力增强，而不是研究范式的真正重排。

如果说机器人学过去数十年的主轴，是“如何让机器稳定地在物理世界中行动”，那么大模型时代带来的新增量，主要在于“如何把更高层的语义理解、世界知识和任务规划能力接入这个行动系统”。这也是为什么近几年具身智能重新升温。语言开始成为更自然的任务接口，多模态基础模型开始为机器人提供更统一的表示层，长上下文与推理能力则让复杂任务分解和跨任务知识迁移变得更可行。

以 \href{https://arxiv.org/abs/2303.03378}{PaLM-E} 为代表的工作提出了“具身多模态语言模型”的思路，也就是把视觉、机器人状态和语言输入统一接入同一大模型体系；\href{https://arxiv.org/abs/2212.06817}{RT-1} 与 \href{https://arxiv.org/abs/2307.15818}{RT-2} 则进一步把机器人控制组织成更可扩展的多任务视觉-语言-动作学习框架。这些工作并不意味着机器人问题已被解决，但它们确实重新定义了机器人感知、语义理解和动作生成之间的接口。

大模型带来的第一个变化，是高层任务接口的统一。过去机器人系统常常依赖硬编码任务模板、状态机和特定场景规则，而大模型让自然语言、开放指令和语义检索更容易成为统一的人机接口。这意味着机器人第一次有机会不再为每个任务单独设计上层入口，而是在更通用的语义空间中接收任务。

第二个变化，是多模态表征成为系统组织中心。传统系统往往把图像、状态、文本和动作分别交给不同模块处理，而大模型路线试图让它们进入同一个表示或至少共享接口体系。这种重组让跨任务知识迁移、长时程上下文利用和不同数据源的联合训练变得更现实，也解释了为什么 VLM 到 VLA 被看作一个自然延伸方向。

第三个变化，是“从会做单任务到会做任务分解”的期待被显著提高。大模型的推理、分解、记忆与交互式纠错能力，使机器人不再只被期待学会一条策略，而被期待能够理解目标、分解步骤、调用技能并在失败后重试。也正因为此，具身智能开始越来越多地与“技能库”“程序化中间表示”“大小脑协同”和“外部记忆”这些概念绑定。

更重要的是，大模型也让机器人领域的难题暴露得更清楚了。语言理解更强，并不自动意味着动作更稳；互联网知识更多，也不自动意味着真实物理交互更可靠；更大的上下文和更长的推理链，也不自动意味着部署时延、安全验证和失效恢复更可控。因此，大模型的价值不在于替代机器人学，而在于迫使整个领域重新思考：哪些能力应该统一，哪些能力必须分层，哪些环节仍需保留强工程约束。

从文献脉络上看，PaLM-E、RT-1 和 RT-2 的重要性，不只在于它们各自性能如何，而在于它们共同推动了一个新的组织框架：机器人不再只是“控制器 + 感知模块”的组合，而是被越来越多地组织成视觉、语言、动作共同进入的具身多模态系统。\href{https://arxiv.org/abs/2303.03378}{PaLM-E} 明确提出了“具身多模态语言模型”这一概念；\href{https://arxiv.org/abs/2212.06817}{RT-1} 说明多任务 Transformer 式控制模型能够在更统一的数据组织下扩展；\href{https://arxiv.org/abs/2307.15818}{RT-2} 则进一步展示了互联网级视觉语言知识如何迁移到机器人任务理解中。它们共同支撑的不是“问题已解”，而是“问题被重新组织”这一判断。

因此，本节的结论不应写成“因为大模型出现，所以具身智能突然诞生”，而应写成：大模型为具身系统提供了新的高层任务接口、表征共享机制与语义迁移能力，同时也让旧问题以更清晰的方式暴露出来。正是在这一意义上，具身智能的重新升温不是一次简单热点循环，而是机器人学主线与大模型能力在同一问题框架中的重新会合。

如果把这种“重新组织”进一步形式化，它更像是把系统主链条中的高层接口改写了：

\begingroup
\scriptsize
\begin{equation*}
\text{classical stack}:\quad o_t \rightarrow \hat x_t \rightarrow \text{planner} \rightarrow \text{controller}
\end{equation*}
\endgroup

\begingroup
\scriptsize
\begin{equation*}
\text{LLM-era stack}:\quad (o_t,\ l,\ m) \rightarrow z_t \rightarrow \text{planner/skill router} \rightarrow \text{controller}
\end{equation*}
\endgroup

差别不在于后者抹掉了控制器，而在于语言 \(l\) 与外部记忆 \(m\) 被正式接入了系统接口，并迫使表征层与规划层承担更多原先由人工规则完成的组织工作。也正因为底层物理闭环并未消失，本报告才会反复强调：大模型改变的是系统组织方式，而不是物理世界的约束定律。
